% Agent 执行循环图
% User Goal → Plan → Tool Call → Observe → 判断 → 完成 / 重规划
\documentclass[border=5pt,tikz]{standalone}
\usepackage{tikz}
\usetikzlibrary{shapes.geometric, arrows.meta, positioning, fit, backgrounds, calc, decorations.pathreplacing}
\usepackage{xeCJK}
\setCJKmainfont{STSong}

\begin{document}
\begin{tikzpicture}[
    % 节点样式
    node distance=0.6cm,
    startstop/.style={rectangle, rounded corners=8pt, draw=green!60!black, fill=green!8,
                      minimum width=2.8cm, minimum height=1.0cm, font=\bfseries\small, align=center},
    process/.style={rectangle, rounded corners=4pt, draw=blue!60!black, fill=blue!6,
                    minimum width=2.4cm, minimum height=1.1cm, font=\footnotesize\bfseries, align=center},
    toolnode/.style={rectangle, rounded corners=4pt, draw=orange!70!black, fill=orange!6,
                     minimum width=2.4cm, minimum height=0.9cm, font=\footnotesize, align=center},
    decision/.style={diamond, draw=red!60!black, fill=red!5, aspect=1.5,
                     minimum width=1.8cm, minimum height=1.0cm, font=\footnotesize\bfseries, align=center},
    arr/.style={-{Stealth[scale=0.9]}, line width=1.2pt, draw=black!50},
    arrback/.style={-{Stealth[scale=0.9]}, line width=1.2pt, draw=red!40, densely dashed},
    annot/.style={font=\tiny, text=black!50},
]

% ==================== 主循环 ====================
\node[startstop] (start) {用户目标\\(User Goal)};

\node[process] (plan) [below=of start] {规划\\(Planning)};

\node[decision] (decide) [below=of plan] {需要\\工具？};

\node[process] (toolcall) [right=0.8cm of decide] {调用工具\\(Tool Call)};

\node[process] (observe) [below=of decide] {观察结果\\(Observe)};

\node[decision] (done) [below=of observe] {目标\\达成？};

\node[startstop, fill=green!20] (finish) [below=of done] {完成\\(Complete)};

% 工具示例
\node[toolnode] (tool1) [above right =of toolcall] {\texttt{read\_file()}\\\footnotesize 读取仿真 log};
\node[toolnode] (tool2) [below=of tool1] {\texttt{search\_kb()}\\\footnotesize 搜索知识库};
\node[toolnode] (tool3) [below=of tool2] {\texttt{run\_sim()}\\\footnotesize 启动仿真};

% ==================== 箭头 ====================
\draw[arr] (start) -- (plan);
\draw[arr] (plan) -- (decide);
\draw[arr] (decide) -- (toolcall) node[midway, above, annot] {是};
\draw[arr] (decide) -- (observe) node[midway, left, annot] {否};
\draw[arr] (toolcall) |- (observe);
\draw[arr] (observe) -- (done);
\draw[arr] (done) -- (finish) node[midway, right, annot] {是};

% 重规划回路
\draw[arrback] (done.west) -- ++(-0.5,0) |- (plan.west)
    node[pos=0, left, font=\tiny\bfseries, text=red!60!black] {否，重规划};

% 工具连接
\draw[->, draw=orange!40, thin, dashed] (toolcall.east) -- (tool1.west);
\draw[->, draw=orange!40, thin, dashed] (toolcall.east) -- (tool2.west);
\draw[->, draw=orange!40, thin, dashed] (toolcall.east) -- (tool3.west);

% ==================== 标签 ====================
% Memory 环
\node[rectangle, draw=purple!40, fill=purple!3, rounded corners=6pt,
    minimum width=1.8cm, minimum height=6.5cm, font=\footnotesize\bfseries, text=purple!70!black,
    align=center] (memory) [left=4.2cm of toolcall] {记忆\\(Memory)};

\draw[<->, draw=purple!30, thin] (memory.east) -- (plan.west);
\draw[<->, draw=purple!30, thin] (memory.east) -- (observe.west);

% ==================== 风险警告 ====================
\node[font=\footnotesize\bfseries, text=red!70!black, align=left]
    (warn) [below=of finish, yshift=0.2cm] {
    ⚠ 风险：规划不可靠 · 工具调用不稳定 · 幻觉在多步中放大\\
    \quad\; 在 EDA 场景中必须\textbf{强约束}：只读自动、写入审核、仿真需确认
};

\end{tikzpicture}
\end{document}
